% Part: first-order-logic
% Chapter: model-theory
% Section: overspill

\documentclass[../../../include/open-logic-section]{subfiles}

\begin{document}

\olfileid{mod}{bas}{ove}
\olsection{అధిప్రసరణ (ఓవర్‌స్పిల్)}

\begin{thm}
\ollabel{overspill}
వాక్యాల సమితి $\Gamma$కు యథేచ్ఛగా పెద్ద పరిమిత
నమూనాలు ఉంటే, దానికి ఒక అనంత నమూనా ఉంటుంది.
\end{thm}

\begin{proof}
$c_0$, $c_1$, \dots అనే లెక్కించదగినన్ని కొత్త స్థిరసంకేతాలను చేర్చి
$\Gamma$ యొక్క భాషను విస్తరించండి; $\Gamma \cup \{c_i \neq c_j :
i \neq j\}$ సమితిని పరిశీలించండి. $\Gamma$కు యథేచ్ఛగా పెద్ద పరిమిత
నమూనాలు ఉన్నాయనడం అంటే, ప్రతి $m>0$కు $n\ge m$ అయ్యేలా
కార్డినాలిటీ~$n$ గల $\Gamma$ నమూనా ఉందని అర్థం. దీనివల్ల $\Gamma
\cup \{c_i \neq c_j : i \neq j\}$ పరిమితంగా సంతృప్తిపరచదగినదని
వస్తుంది. సంహతత్వం వల్ల $\Gamma \cup \{c_i \neq c_j : i \neq j\}$కు
ఒక నమూనా $\Struct M$ ఉంటుంది. అది అన్ని $c_i \neq c_j$ అసమానతలను
సంతృప్తిపరుస్తుంది కాబట్టి దాని వ్యక్తి క్షేత్రం అనంతమై ఉండాలి.
\end{proof}

\begin{prop}
\ollabel{inf-not-fo}
ఏ మొదటిస్థాయి భాషలోనూ, \tetoken{ఒక నిర్మాణం}{!!a{structure}}~$\Struct M$లో
సత్యమవడం మరియు ఆ \tetoken{నిర్మాణపు}{!!{structure}} వ్యక్తి క్షేత్రం
$\Domain{M}$ అనంతమవడం పరస్పరం తుల్యమయ్యే వాక్యం~$!A$ లేదు.
\end{prop}

\begin{proof}
అటువంటి $!A$ ఉందనుకుంటే, దాని నిషేధం $\lnot !A$ అన్ని పరిమిత
\tetoken{నిర్మాణాల్లో}{!!{structure}s}, వాటిలో మాత్రమే సత్యమవుతుంది. కాబట్టి
ఆ నిషేధానికి యథేచ్ఛగా పెద్ద పరిమిత నమూనాలు ఉంటాయి, కానీ అనంత నమూనా
ఉండదు; ఇది \olref{overspill}కు విరుద్ధం.
\end{proof}

\end{document}
